\documentclass[11pt]{article}
\usepackage[margin=1.15in]{geometry}
\usepackage{amsmath,amssymb,amsthm,mathtools}
\usepackage[T1]{fontenc}
\usepackage{lmodern}
\usepackage{booktabs}
\usepackage{enumitem}
\usepackage[colorlinks=true,linkcolor=blue!40!black,citecolor=blue!40!black,urlcolor=blue!40!black]{hyperref}
\theoremstyle{plain}
\newtheorem{theorem}{Theorem}[section]
\newtheorem{proposition}[theorem]{Proposition}
\newtheorem{lemma}[theorem]{Lemma}
\newtheorem{corollary}[theorem]{Corollary}
\theoremstyle{definition}
\newtheorem{definition}[theorem]{Definition}
\newtheorem{conjecture}[theorem]{Conjecture}
\newtheorem{question}[theorem]{Question}
\theoremstyle{remark}
\newtheorem{remark}[theorem]{Remark}
\newcommand{\R}{\mathbb{R}}
\newcommand{\N}{\mathbb{N}}
\newcommand{\Q}{\mathbb{Q}}
\newcommand{\C}{\mathbb{C}}
\newcommand{\Gr}{\mathrm{Gr}}
\newcommand{\dd}{\mathrm{d}}
\newcommand{\eps}{\varepsilon}
\DeclareMathOperator{\sgn}{sgn}

\title{Answers to some recent conjectures and questions in classical analysis, with computer assistance}
\author{Jacob Goodchild\thanks{Directing author. The mathematical work was carried out by an AI system, Claude (Anthropic), via Claude Code; see the statement on the use of AI at the end of Section~\ref{sec:intro}.}}
\date{November 2026}
\begin{document}
\maketitle
\begin{abstract}
We answer, or make progress on, several questions and conjectures posed in recent arXiv papers in classical analysis. We prove two conjectures of Fan on Riesz capacity of three-point sets on the circle and of vertex sets of regular odd polygons. We show that a conjecture of Xu and Zhao on the sign thresholds of higher-order Gregory coefficients is false and compute these thresholds rigorously for orders up to $10$. We prove by computer-assisted ball arithmetic that the reciprocal of the diagonal Pad\'e denominator of $e^z$ has positive Taylor coefficients for every odd degree $3\le m\le361$, a case left open in the literature, and describe numerically where positivity first fails for fixed $m$. We also give explicit values for series of products of consecutive harmonic numbers, short answers to questions of Finch on Abel's equation and iterates, two of three open problems of Wi\'sniewska on Gurland's ratio, and numerical values of two completely monotonic degrees. Theorems, computer-assisted proofs and numerical results are clearly separated. The work was carried out with extensive use of an AI system (see the statement on the use of AI).
\end{abstract}
\section{Introduction}\label{sec:intro}

This paper collects answers to a number of specific questions and conjectures that appeared in recent arXiv papers (2021--2026) in classical analysis: inequalities and sharp constants, series and sums, special functions, iteration and asymptotics. The questions were chosen from the literature by an AI system, by the rule that someone had written the question down (as a conjecture, an open problem, or a request for a value) and that it could be attacked by calculus together with computer algebra, interval arithmetic or exact computation. For each question we state the original wording and its source, we state our result, and we give the proof or, for numerical results, the method and the precision. We distinguish three kinds of statements throughout.

\begin{itemize}[leftmargin=2em]
\item \textbf{Theorems}: complete proofs in the text (Section~\ref{sec:riesz}; Propositions~\ref{prop:delta} and~\ref{prop:gurland}, which are short and elementary).
\item \textbf{Computer-assisted theorems}: statements whose proof consists of a written argument plus a finite computation in exact or interval arithmetic that we describe precisely (Theorems~\ref{thm:gregory}(ii) and~\ref{thm:pade}, and the exact reductions in Section~\ref{sec:harmonic}).
\item \textbf{Numerical results}: values or laws obtained by high-precision computation and checked by independent code, but not proved (for example the tables in Sections~\ref{sec:finch} and~\ref{sec:stirling} and the law for $k^*(m)$ in Section~\ref{sec:kstar}). These are never presented as theorems.
\end{itemize}

\paragraph{Summary of results.}
\begin{enumerate}[leftmargin=2em]
\item (Section~\ref{sec:riesz}) We prove two conjectures of Fan~\cite{Fan} on Riesz capacity of three-point sets on a circle and of the vertex sets of regular odd polygons (Theorem~\ref{thm:riesz}).
\item (Section~\ref{sec:gregory}) We show that a conjecture of Xu and Zhao~\cite{XZ} on the sign thresholds $\Gr(\ell)$ of higher-order Gregory coefficients is false, compute $\Gr(\ell)$ rigorously for $\ell\le10$ (and numerically up to $14$), and describe the observed growth.
\item (Section~\ref{sec:pade}) For the denominators $q_{m,m}$ of diagonal Pad\'e approximants to $e^z$ we prove, by computer assistance, that $1/q_{m,m}$ has positive Taylor coefficients for every odd $3\le m\le361$, a case left open in~\cite{Scardino} of a conjecture of Fasi and Higham~\cite{FH}, and we describe numerically where positivity first fails for fixed $m$.
\item (Section~\ref{sec:harmonic}) We give an exact algorithm and explicit values for the series $S_q=\sum_{k\ge1}H_kH_{k+1}\cdots H_{k+q}/(k(k+1)\cdots(k+q))$ studied in~\cite{CN}.
\item (Section~\ref{sec:finch}) We give a derivation of Finch's conjectured values of the constant between two solutions of Abel's equation (an easy consequence of classical Fatou-coordinate asymptotics), an exact series and the value of a second derivative asked for in~\cite{FinchIII}, and half-iterate values asked for in~\cite{FinchIV}.
\item (Section~\ref{sec:gurland}) We answer two of the three open problems of Wi\'sniewska~\cite{Wis} on Gurland's ratio (by an elementary argument) and give numerical information on the third.
\item (Section~\ref{sec:stirling}) We give numerical values of two completely monotonic degrees whose exact values are said to be open in~\cite{Kr}.
\end{enumerate}

\paragraph{Results that turned out to be known.} We also found proofs of three statements that, as our literature search showed, had already been settled, and which we therefore do not claim: Conjecture 7.1 of Habib and Hildebrand~\cite{HH} (proved by the same authors a few days later in~\cite{HH2}); Conjecture 4.5 of Chavan, Kobayashi and Layja~\cite{CKL} on the multiple $t$-values $t(\{2\}^N,1)$ (proved by Charlton~\cite{Charlton}); and the two conjectured identities between Abel functions in~\cite{FinchAB} (explained by the author himself in the corrigendum of that paper). A numerical refutation of a Mahler-measure conjecture and its explanation, which we also reproduced, are in~\cite{Zheng}. We mention them to give an accurate account of what the search found.

\paragraph{Limits of the literature search.} Novelty statements are based on web and arXiv searches carried out in November 2026; citation-count tools were mostly unavailable, so a later or less visible reference to a result may have been missed. ``Not found'' means only that. Several papers cited here appeared in 2026; they are cited by arXiv number and were read in their arXiv versions.

\paragraph{Code availability.} All code, logs and verification scripts are in the repository \url{https://github.com/JacobGoodchild/findformula3}, under \texttt{results/}; the file \texttt{VERIFICATION.md} lists, for each result, what was checked and by which method.

\paragraph{Statement on the use of AI.}
This work was carried out with extensive use of an AI system, Claude (Anthropic), via the Claude Code software, under the direction of Jacob Goodchild. The AI system chose the problems from the recent literature, carried out the derivations, proofs and computations, searched and checked the literature, and wrote the manuscript. Proofs were checked by independent AI referee passes and by computer algebra and interval arithmetic, but have not been checked by a human expert or peer reviewed. Novelty assessments are based on literature searches and may be incomplete. Jacob Goodchild directed the work and is responsible for its publication.

\section{Riesz energy of three points on a circle}\label{sec:riesz}

\subsection{The question}
Let $K\subset\R^2$ be compact and $r>0$. Write
\[
U_r(K)\;=\;\max_{\mu}\iint |x-y|^{r}\,\dd\mu(x)\,\dd\mu(y),
\]
the maximum over Borel probability measures $\mu$ on $K$ (a maximiser is called an \emph{equilibrium measure}). In the notation of Fan~\cite{Fan}, $U_r(K)=V_{-r}(K)$ and the Riesz capacity is $\mathrm{Cap}_{-r}(K)=U_r(K)^{1/r}$, so maximising capacity and maximising $U_r$ are the same problem. Fan studies the vertex set $P_N$ of the regular $N$-gon inscribed in the unit circle and states two conjectures, which we quote.

\begin{conjecture}[Fan~\cite{Fan}, Conjecture 5.9]\label{conj:fan59}
Suppose $r\ge2$. Let $A,B,C$ be three points on the unit circle. Fix $B$ and $C$ such that they are symmetric with respect to the $x$-axis, having polar angles $\varphi$ and $-\varphi$, respectively, where $\varphi\in(\pi/2,2\pi/3]$. Move $A$ along the circle with polar angle $\psi\in[0,2\pi-3\varphi]$, so that $BC$ remains the longest side of the triangle $\triangle ABC$. Then the capacity of the three-point set $\{A,B,C\}$ is maximized when $\psi=2\pi-3\varphi$, i.e., when $AC=BC$.
\end{conjecture}

\begin{conjecture}[Fan~\cite{Fan}, Conjecture 5.6]\label{conj:fan56}
Let $r>2$ and let $P_N$ be the vertex set of a regular $N$-gon inscribed in the unit circle ($N\ge3$) with $N$ odd. Let $\theta_N=\frac{(N-1)\pi}{2N}$. Then the energy equals $U_r(P_N)=\dfrac{2^{r-1}\sin^r\theta_N}{1-2^{r-2}\cos^r\theta_N}$, and an equilibrium measure is supported on the three points $1$, $e^{2\theta_N i}$, $e^{2(\pi-\theta_N)i}$ (up to rotation), with masses $\dfrac{2b^r-a^r}{4b^r-a^r},\ \dfrac{b^r}{4b^r-a^r},\ \dfrac{b^r}{4b^r-a^r}$ respectively, where $a=2\sin(\pi/N)$ and $b=2\sin\theta_N$.
\end{conjecture}

Fan proves the special case of Conjecture~\ref{conj:fan59} in which $A$ is replaced by the symmetric configuration (her Lemma 5.8, restated below), shows that Conjecture~\ref{conj:fan59} together with that lemma implies Conjecture~\ref{conj:fan56}, and describes her attempt to prove Conjecture~\ref{conj:fan59} as unsuccessful. We prove both conjectures. The proof is a calculus argument; its only external inputs are three results quoted in~\cite{Fan}, which we use as stated there.

\begin{theorem}\label{thm:riesz}
Conjecture~\ref{conj:fan59} and Conjecture~\ref{conj:fan56} are true. More precisely, in the setting of Conjecture~\ref{conj:fan59}, the energy $U_r(\{A,B,C\})$ is a nondecreasing function of $\psi\in[0,2\pi-3\varphi]$.
\end{theorem}

\noindent\emph{Notation.} In Conjecture~\ref{conj:fan56} the symbols $a,b$ are Fan's; in the proof below they correspond to $\beta$ and $\ell$. Elsewhere $a,b,c$ denote side lengths of a triangle.\medskip

\subsection{Inputs taken from the literature}
\begin{itemize}[leftmargin=2em]
\item[\textbf{(I1)}] \emph{Energy of three points} (Clark--Laugesen~\cite{CL}, Theorem 8, quoted as Lemma 5.7 in~\cite{Fan}). Let $T\subset\R^2$ have three points with mutual distances $a,b,c$, $0<a,b\le c$, and let $r>0$. If $a^r+b^r\le c^r$ then $U_r(T)=c^r/2$; otherwise
\[
U_r(T)=\frac{2a^rb^rc^r}{4a^rb^r-(a^r+b^r-c^r)^2}.
\]
\item[\textbf{(I2)}] \emph{Support} (Bj\"orck~\cite{Bj}, Theorem 12, quoted in~\cite{Fan} before Lemma 5.2). If $r>2$, every compact $K\subset\R^2$ has an equilibrium measure supported on at most three points, which are extreme points of the convex hull of $K$. In particular $U_r(P_N)=\max\{U_r(T): T\subset P_N,\ |T|\le3\}$.
\item[\textbf{(I3)}] \emph{Isosceles triples} (Fan, Lemma 5.8). Let $r\ge2$, $A=(1,0)$, $B=e^{i\alpha}$, $C=e^{-i\alpha}$. Then $U_r(\{A,B,C\})$ is nondecreasing in $\alpha\in[2\pi/3,\pi)$.
\end{itemize}
We did not open the papers of Bj\"orck and Clark--Laugesen; (I1) and (I2) are used exactly as they are stated in~\cite{Fan}.

\subsection{Notation and a formula for the derivative}
Throughout, $r\ge 2$. For a triple with sides $a,b\le c$ put $x=a^r$, $y=b^r$, $z=c^r$ and $s=x+y-z$. If $s\le0$ we are on the \emph{two-point branch}, $U=z/2$; if $s>0$ we are on the \emph{three-point branch}, and then, since $z\ge\max(x,y)$ gives $s\le\min(x,y)$,
\[
D:=4xy-s^2\ \ge\ 4xy-\min(x,y)^2\ >0,\qquad U=\frac{2xyz}{D},\qquad U-\frac z2=\frac{z\,s^2}{2D}>0 .
\]
At $s=0$ both expressions equal $z/2$, so $U$ is a continuous function of $(x,y,z)$.

Fix $\varphi\in(\pi/2,2\pi/3]$ and let $B=e^{i\varphi}$, $C=e^{-i\varphi}$, $A=e^{i\psi}$ with $0\le\psi\le 2\pi-3\varphi$. Put
\[
\lambda=\frac{\varphi-\psi}{2},\qquad \eta=\frac{\varphi+\psi}{2},\qquad\text{so}\qquad 0\le\lambda\le\eta\le\pi-\varphi<\frac\pi2<\varphi .
\]
(The inequality $\eta\le\pi-\varphi$ is $\psi\le2\pi-3\varphi$; $\lambda>0$ for $\psi<\varphi$, which holds because $2\pi-3\varphi<\varphi$ when $\varphi>\pi/2$.) The chord lengths are
\[
a:=|AB|=2\sin\lambda,\qquad b:=|AC|=2\sin\eta,\qquad c:=|BC|=2\sin\varphi .
\]
Since $\sin$ is increasing on $[0,\pi/2]$ and $\sin\varphi=\sin(\pi-\varphi)$, we have $a\le b\le c$, with $b=c$ exactly at $\psi=2\pi-3\varphi$ and $a=b$ exactly at $\psi=0$. So $BC$ is the longest side, as in Conjecture~\ref{conj:fan59}. As $\psi$ increases, $a$ decreases and $b$ increases, $c$ is constant. Write $F(\psi)=U_r(\{A,B,C\})$; $F$ is continuous.

\begin{lemma}[derivative]\label{lem:deriv}
On a $\psi$-interval where $s>0$ and $0<\psi<2\pi-3\varphi$ (so $0<\lambda<\mu$ and $b<c$),
\[
\frac{\dd\log F}{\dd\psi}=\frac{s\,r}{2D}\Big[(z+y-x)\cot\eta-(z+x-y)\cot\lambda\Big].
\]
\end{lemma}
\begin{proof}
From $\log U=\log2+\log x+\log y+\log z-\log D$ and $D=4xy-s^2$, $s=x+y-z$,
\[
\frac{\partial\log U}{\partial x}=\frac1x-\frac{4y-2s}{D}=\frac{D-x(4y-2s)}{xD}=\frac{s(2x-s)}{xD}=\frac{s\,(z+x-y)}{xD},
\]
and symmetrically $\partial_y\log U=s(z+y-x)/(yD)$. Now $\dd\lambda/\dd\psi=-\tfrac12$, $\dd\eta/\dd\psi=\tfrac12$, so $\dd\log x/\dd\psi=-\tfrac r2\cot\lambda$ and $\dd\log y/\dd\psi=\tfrac r2\cot\eta$, while $z$ does not depend on $\psi$. The chain rule gives the formula.
\end{proof}

\subsection{Monotonicity of $F$}
\begin{lemma}[an elementary inequality]\label{lem:ineq}
Let $0<p<q<1$, let $r\ge2$ and suppose $p^r+q^r>1$. Then $q^r-p^r\ \ge\ q^2-p^2$.
\end{lemma}
\begin{proof}
Put $G(\rho)=q^\rho-p^\rho-(q^2-p^2)$ for real $\rho\ge2$, so $G(2)=0$ and $G'(\rho)=q^\rho\ln q-p^\rho\ln p$. Since $\ln p,\ln q<0$, $G'(\rho)\ge0$ is equivalent to $(q/p)^\rho\le \ln p/\ln q$. The left side is strictly increasing in $\rho$ (because $q>p$), so the set $\{\rho\ge2:G'(\rho)\ge0\}$ is an initial segment of $[2,\infty)$ (possibly empty): $G$ is nondecreasing on it and strictly decreasing afterwards. In either case, on every interval $[2,\rho_1]$ we have $G\ge\min\{G(2),G(\rho_1)\}$.

The function $\rho\mapsto p^\rho+q^\rho$ is strictly decreasing and tends to $0$ as $\rho\to\infty$; let $\rho_1$ be the root of $p^{\rho_1}+q^{\rho_1}=1$. Because $p^r+q^r>1$ we have $2\le r<\rho_1$, so it suffices to prove $G(\rho_1)\ge0$.

Put $w=q^{\rho_1}$. Then $p^{\rho_1}=1-w$, $w\in(\tfrac12,1)$ (as $q>p$), and with $\sigma=2/\rho_1\in(0,1]$ we have $q^2=w^\sigma$, $p^2=(1-w)^\sigma$. Thus
\[
G(\rho_1)=(2w-1)-\big(w^\sigma-(1-w)^\sigma\big)=\chi(1)-\chi(\sigma),\qquad \chi(\sigma):=w^\sigma-(1-w)^\sigma ,
\]
and it suffices to show that $\chi$ is nondecreasing on $(0,1]$. We have $\chi'(\sigma)=(1-w)^\sigma|\ln(1-w)|-w^\sigma|\ln w|$, so
\[
\frac{\chi'(\sigma)}{w^\sigma|\ln w|}=\Big(\frac{1-w}{w}\Big)^{\sigma}\frac{|\ln(1-w)|}{|\ln w|}-1\ \ge\ \frac{(1-w)|\ln(1-w)|}{w|\ln w|}-1=\frac{k(1-w)}{k(w)}-1,
\]
where $k(u)=-u\ln u$ and we used $\frac{1-w}{w}<1$, $\sigma\le1$. It remains to see $k(1-w)\ge k(w)$ for $w\in[\tfrac12,1)$. Put $u=1-w\in(0,\tfrac12]$ and $E(u)=k(u)-k(1-u)$. Then $E$ is continuous on $[0,\tfrac12]$ with $E(0)=E(\tfrac12)=0$, and $E''(u)=-\frac1u+\frac1{1-u}=\frac{2u-1}{u(1-u)}<0$ on $(0,\tfrac12)$, so $E$ is concave and $E\ge0$. Hence $\chi'\ge0$.
\end{proof}

\begin{proposition}\label{prop:mono}
$F$ is nondecreasing on $[0,2\pi-3\varphi]$.
\end{proposition}
\begin{proof}
On the open set where $s>0$ and $0<\psi<2\pi-3\varphi$ we use Lemma~\ref{lem:deriv}. Here $s>0$, $D>0$, $\cot\eta>0$ (as $\eta<\pi/2$), $\cot\lambda>0$, and $z+x-y\ge x>0$ (as $y\le z$). So the sign of $\dd\log F/\dd\psi$ is the sign of
\[
R_r-\frac{\cot\lambda}{\cot\eta},\qquad R_r:=\frac{z+y-x}{z+x-y}.
\]
Put $p=a/c$, $q=b/c$ (so $0<p<q<1$) and $t_r=q^r-p^r\in(-1,1)$. Then $R_r=\dfrac{1+t_r}{1-t_r}$, which is an increasing function of $t_r$. For $r=2$ we compute with $\varphi=\lambda+\eta$ and $\sin^2\eta-\sin^2\lambda=\sin(\eta-\lambda)\sin(\eta+\lambda)=\sin\psi\sin\varphi$:
\[
R_2=\frac{\sin^2\varphi+\sin\varphi\sin\psi}{\sin^2\varphi-\sin\varphi\sin\psi}=\frac{\sin\varphi+\sin\psi}{\sin\varphi-\sin\psi}=\frac{2\sin\eta\cos\lambda}{2\cos\eta\sin\lambda}=\frac{\cot\lambda}{\cot\eta},
\]
where $\sin\varphi+\sin\psi=2\sin\eta\cos\lambda$ and $\sin\varphi-\sin\psi=2\cos\eta\sin\lambda$. Hence the sign of $\dd\log F/\dd\psi$ is the sign of $R_r-R_2$, i.e.\ of $t_r-t_2$, i.e.\ of $q^r-p^r-(q^2-p^2)$.

We check the hypotheses of Lemma~\ref{lem:ineq}: $0<p<q<1$ was noted above, and $p^r+q^r>1$ is exactly $s>0$. By Lemma~\ref{lem:ineq}, $t_r\ge t_2$, hence $\dd\log F/\dd\psi\ge0$ wherever $s>0$ in the open interval.

It remains to treat the branches. Suppose $s(\psi_1)>0$ for some $\psi_1\in[0,2\pi-3\varphi)$, and let $\psi_2=\sup\{\psi'\in(\psi_1,2\pi-3\varphi]: s>0\text{ on }[\psi_1,\psi')\}$. On $[\psi_1,\psi_2)$ the function $F$ is continuous and nondecreasing (it is differentiable with nonnegative derivative on the open set where $s>0$ and $0<\psi<2\pi-3\varphi$), so $F(\psi_2)\ge F(\psi_1)>z/2$ by continuity. If $\psi_2<2\pi-3\varphi$ then $s(\psi_2)=0$ and $F(\psi_2)=z/2$, a contradiction; hence $\psi_2=2\pi-3\varphi$, i.e.\ once $s>0$ it stays positive up to the endpoint, and $F$ is nondecreasing on $[\psi_1,2\pi-3\varphi]$. On the set where $s\le0$, $F\equiv z/2$ is constant, and $s>0$ at $\psi$ implies $s>0$ later, so this set is an initial interval $[0,\psi_*]$ (possibly empty) on which $F=z/2$. Since $F$ is continuous and $F\ge z/2$ always, $F$ is nondecreasing on $[0,2\pi-3\varphi]$. The degenerate case $\varphi=2\pi/3$ (where $2\pi-3\varphi=0$) is trivial.
\end{proof}

\begin{proof}[Proof of Conjecture~\ref{conj:fan59}]
By Proposition~\ref{prop:mono}, $F(\psi)\le F(2\pi-3\varphi)$, i.e.\ the energy, and hence the capacity $U_r^{1/r}$, is maximal at $\psi=2\pi-3\varphi$, where $b=c$ (that is $AC=BC$).
\end{proof}

\begin{remark}
The inequality of Lemma~\ref{lem:ineq} gives only $\ge$, so we claim monotonicity (not strict monotonicity) of $F$, which is all that Conjecture~\ref{conj:fan59} asserts. At $r=2$, $t_2-t_2=0$ and $F$ is constant, as in Theorem 5.4 of~\cite{Fan}.
\end{remark}

\subsection{Proof of Conjecture~\ref{conj:fan56}}
Let $N\ge3$ be odd, $r>2$, $\theta=\theta_N=\frac\pi2-\frac\pi{2N}$, and let $v_k=e^{2\pi ik/N}$. Put $\ell=2\sin\theta$ (the longest chord of $P_N$) and $\beta=2\sin(\pi/N)$. Since $\pi/N=\pi-2\theta$ we have $\beta=2\sin2\theta=4\sin\theta\cos\theta$.

\emph{The conjectured triple.} Let $T_N=\{1,\,e^{i(\pi-\pi/N)},\,e^{i(\pi+\pi/N)}\}$; this is the triple of Conjecture~\ref{conj:fan56} because $2\theta_N=\pi-\pi/N$ and $2(\pi-\theta_N)=\pi+\pi/N$. For odd $N$ these are vertices of $P_N$ ($\frac{N\mp1}{2}$ are integers). $T_N$ is isosceles with two sides $\ell$ and base $\beta\le\ell$ (this is $\sin\frac{\pi}{2N}\le\frac12$, true for $N\ge3$). Applying (I1) with $c=\ell$ and the two other sides $\ell$, $\beta$ (so $\ell^r+\beta^r>\ell^r$: three-point branch),
\[
U_r(T_N)=\frac{2\ell^{2r}\beta^r}{4\ell^r\beta^r-\beta^{2r}}=\frac{2\ell^{2r}}{4\ell^r-\beta^r}=\frac{2^{r-1}\sin^r\theta}{1-2^{r-2}\cos^r\theta},
\]
where the last equality uses $\ell=2\sin\theta$, $\beta=4\sin\theta\cos\theta$. This is the energy claimed in Conjecture~\ref{conj:fan56}. Moreover let $\mu_N$ be the measure on $T_N$ with the masses of Conjecture~\ref{conj:fan56} (the apex $1$ carries $\frac{2\ell^r-\beta^r}{4\ell^r-\beta^r}$, and the other two points carry $\frac{\ell^r}{4\ell^r-\beta^r}$ each; here $a=\beta$, $b=\ell$ in the notation of the conjecture). These masses are positive ($2\ell^r>\beta^r$) and sum to $1$, and a direct computation gives
\[
\iint|x-y|^r\dd\mu_N\dd\mu_N=2\Big[2\,m_1m_2\ell^r+m_2^2\beta^r\Big]=\frac{2\ell^{2r}\,(4\ell^r-\beta^r)}{(4\ell^r-\beta^r)^2}=U_r(T_N).
\]

\emph{Upper bound.} By (I2) it suffices to show $U_r(T)\le U_r(T_N)$ for every $T\subset P_N$ with at most three points. Single points give $U_r=0$, and for a pair at distance $d$ one has $U_r=\max_{0\le p\le1}2p(1-p)d^r=d^r/2$, so pairs give $U_r=d^r/2\le\ell^r/2\le U_r(T_N)$ (where $d\le\ell$ is the distance, and $U_r(T_N)\ge\ell^r/2$ because $U\ge z/2$ in (I1)). Let $T=\{A,B,C\}\subset P_N$ be a triple with longest side $BC$.
\begin{itemize}[leftmargin=2em]
\item If $a^r+b^r\le c^r$ (two-point branch), $U_r(T)=c^r/2\le\ell^r/2\le U_r(T_N)$.
\item If $a^r+b^r>c^r$, then, since $(a/c)^r\le(a/c)^2$ and $(b/c)^r\le (b/c)^2$ for $r>2$, we get $a^2+b^2>c^2$, so all angles of the triangle are acute and $T$ is an acute inscribed triangle: its three arcs are all smaller than $\pi$. Rotate so that the midpoint of the arc $BC$ not containing $A$ is $e^{i\pi}$; if that arc has length $2\pi-2\varphi$ (which is $<\pi$, so $\varphi>\pi/2$) then $B,C=e^{\pm i\varphi}$ (after naming), and $A=e^{i\psi}$ lies on the complementary arc, $|\psi|<\varphi$; after a reflection in the $x$-axis, $\psi\in[0,\varphi)$; the arcs $AB,AC$ have lengths $\varphi-\psi,\ \varphi+\psi$; $BC$ longest means $2\pi-2\varphi\ge\varphi+\psi$, i.e.\ $\psi\le2\pi-3\varphi$; and since the three arcs sum to $2\pi$ the longest is at least $2\pi/3$, so $\varphi\le2\pi/3$. Thus $T$ is exactly in the configuration of Theorem~\ref{thm:riesz} and, by Proposition~\ref{prop:mono}, $U_r(T)\le U_r(T')$ where $T'$ is the configuration with $\psi=2\pi-3\varphi$, i.e.\ $|AC|=|BC|=:L=2\sin\varphi$. An isosceles triangle inscribed in the unit circle with legs $L$ is unique up to congruence (its apex is a point of the circle and the other two vertices are the two points of the circle at distance $L$ from it), hence $T'$ is congruent to $\{(1,0),e^{i\alpha},e^{-i\alpha}\}$ with $2\sin(\alpha/2)=L$, i.e.\ $\alpha=2\pi-2\varphi\in[2\pi/3,\pi)$ (as $\varphi\in(\pi/2,2\pi/3]$). As $L=|BC|$ is a chord of $P_N$, $L\le\ell=2\sin\theta$, so $\alpha/2\le\theta$ (both lie in $(0,\pi/2)$) and $\alpha\le2\theta=\pi-\pi/N$, where $2\theta\ge2\pi/3$ because $N\ge3$. By (I3), $U_r(T')\le U_r(\{(1,0),e^{i(\pi-\pi/N)},e^{-i(\pi-\pi/N)}\})=U_r(T_N)$.
\end{itemize}
Hence $U_r(P_N)\le U_r(T_N)=\iint|x-y|^r\dd\mu_N\dd\mu_N\le U_r(P_N)$, so $U_r(P_N)=U_r(T_N)$ and $\mu_N$ is an equilibrium measure of $P_N$, supported on the three points of Conjecture~\ref{conj:fan56}. This proves Conjecture~\ref{conj:fan56}. \hfill$\square$

\begin{remark}
We prove existence of an equilibrium measure with the stated support and masses; we do not claim uniqueness up to rotation (Fan notes the same limitation of her numerical evidence). When $3\mid N$ the equilateral triangle of vertices is an admissible acute triple covered by the case analysis above.
\end{remark}

\subsection{Checks}
The derivative formula, Lemma~\ref{lem:ineq}, and the assembly of $U_r(T_N)$ were checked with computer algebra and high-precision numerics, and the conclusion of Conjecture~\ref{conj:fan56} was compared with a brute-force maximisation over all vertex triples for odd $N\le31$ and several values of $r>2$ (see the repository, \texttt{results/riesz\_triangle}). These checks are not part of the proof.

\section{Sign thresholds of higher-order Gregory coefficients}\label{sec:gregory}

\subsection{The question}
Let $G_n^{(\ell)}$ be the Taylor coefficients of $\big(x/\ln(1+x)\big)^\ell=\sum_{n\ge0}G_n^{(\ell)}x^n$ (higher-order Gregory coefficients; $\ell=1$ gives the classical Gregory coefficients). Xu and Zhao~\cite{XZ} prove (their Theorem 5.1) that for each $\ell\ge1$ there is a positive integer $\Gr(\ell)$ such that $(-1)^{n-1}G_n^{(\ell)}>0$ for all $n\ge\Gr(\ell)$. Their wording only asserts the existence of \emph{some} integer; they then quote ``exact values'' $\Gr(3)=11$, $\Gr(4)=36$, $\Gr(5)=113$, $\Gr(6)=346$, $\Gr(7)=1030$, so we read $\Gr(\ell)$ as the \emph{least} such integer, i.e.\ $1+{}$the largest $n\ge1$ with $(-1)^{n-1}G_n^{(\ell)}\le0$. (A few small-$n$ coefficients vanish, e.g.\ $G^{(\ell)}_{\ell+1}=0$ for even $\ell$ and $G_\ell^{(\ell)}=0$ for odd $\ell\ge3$; they lie far below the last bad $n$.) They ask:

\begin{conjecture}[Xu--Zhao~\cite{XZ}, Conjecture 6.1]\label{conj:xz}
For all $\ell\ge5$ we have $0.47\cdot3^\ell<\Gr(\ell)<0.477\cdot3^\ell$.
\end{conjecture}

In the same passage the authors remark that the largest real root $r(\ell)$ of the function appearing in their proof ``seems to increase by roughly tripling as $\ell$ increases, similarly to $\{\Gr(\ell)\}$'', and ask whether there is a theoretical reason.

\subsection{Results}
\begin{theorem}\label{thm:gregory}
Conjecture~\ref{conj:xz} is false. Indeed:
\begin{enumerate}[label=(\roman*),leftmargin=2em]
\item $\Gr(5)=113<114.21=0.47\cdot3^5$. This follows from the value $113$ stated in~\cite{XZ}; we also reproduce it.
\item (computer-assisted, rigorous) $\Gr(\ell)$ equals $11,36,113,346,1030,2995,8580,24303$ for $\ell=3,\dots,10$. Thus $\Gr(\ell)<0.47\cdot3^\ell$ also for $\ell=8,9,10$ ($3083.7$, $9251.0$, $27753.0$ against $2995$, $8580$, $24303$). The values for $\ell=3,\dots,7$ agree with~\cite{XZ}; those for $\ell=8,9,10$ are new.
\end{enumerate}
\end{theorem}

\noindent\emph{Proof of (i).} A one-line comparison of the stated value with the stated lower bound.

\noindent\emph{Proof of (ii) (computer-assisted).} The argument has two parts, as in the proof of Theorem 5.1 of~\cite{XZ}.

\emph{Finite range.} The coefficients $G_n^{(\ell)}$, $n\le N_{\max}$, were computed as ball-valued power series (Arb through python-flint, via $\exp(-\ell\log(\log(1+x)/x))$) with every sign certified; the working precision was raised until all coefficients with $n\ge50$ had a definite sign (128 bits sufficed in the independent run below). This gives the largest $n\le N_{\max}$ with the wrong sign.

\emph{Tail.} For $n>\ell+3$, \cite{XZ} writes $(-1)^{n-1}G_n^{(\ell)}=\int_1^\infty h(z)f_n(z)\,\dd z$ (their identity (6)) with explicit $h$, $f_n$, and, by cutting the integral at the largest real root $r=r(\ell)$ of $h$ and bounding the pieces, shows that $(-1)^{n-1}G^{(\ell)}_n>0$ for all $n\ge n_2(\ell)$, where $n_2(\ell)$ is the least $n$ for which an explicit inequality involving integrals over $[1,r]$ and $[0,1/r]$ holds (an inequality which is monotone in $n$). Evaluating that inequality with rigorous interval integration (Arb), with certified isolation of the $\ell-1$ real roots of the polynomial $H_\ell$ of~\cite{XZ}, gives $n_2(\ell)=15,52,168,519,1548,4514,12944,36685$ for $\ell=3,\dots,10$, with positive certified margins; for $\ell=4,\dots,10$ these are the values of Table 2 of~\cite{XZ} ($\ell=3$: $15$ instead of their $16$). All are below the ranges scanned ($N_{\max}\ge6000$ for $\ell\le8$ and $\ge45000$ for $\ell=9,10$). The values of $\Gr(\ell)$ then follow. \hfill$\square$

\smallskip
We stress what this proof takes from~\cite{XZ}: the integral identity (6), the structure of the inequality and its monotonicity in $n$ (we re-checked (6) numerically against exact coefficients and the root counts by certified isolation, but did not re-prove them). The sign scan and the tail evaluation were carried out twice, by independent implementations (the second written without access to the first), with the same values of $\Gr(\ell)$ and, for $\ell=5,\dots,10$, the same thresholds $n_2(\ell)$ (our first, non-rigorous implementation gave $18,53$ instead of $15$/$16$, $52$ for $\ell=3,4$; this does not affect $\Gr(3),\Gr(4)$).

\begin{remark}
Table~1 of~\cite{XZ} lists $r(9)\approx5065$; with $r(\ell)=\exp(\pi\cot(\pi/\ell))$ (see below) one gets $r(9)=5605.2$, so we believe $5065$ is a transposition of digits (their $n_2(9)$ agrees with ours).
\end{remark}

\subsection{More values and the growth of $\Gr(\ell)$}
The same scan, with the tail inequality evaluated numerically (not rigorously, and without certified margins) gives $\Gr(11)=68283$, $\Gr(12)=190698$, $\Gr(13)=530152$, $\Gr(14)=1468691$; an independent sign scan (to $N=1.2\cdot10^5$) reproduced $\Gr(11)$. Table~\ref{tab:gr} lists the data together with $r(\ell)=\exp(\pi\cot(\pi/\ell))$.

\begin{table}[h]
\centering
\caption{$\Gr(\ell)$ for $3\le\ell\le14$ (rigorous for $\ell\le10$; numerical for $\ell=11,\dots,14$). $r(\ell)=\exp(\pi\cot(\pi/\ell))$.}\label{tab:gr}
\begin{tabular}{rrrrrr}
\toprule
$\ell$ & $\Gr(\ell)$ & $\Gr(\ell)/3^\ell$ & $r(\ell)$ & $\Gr(\ell)/r(\ell)$ & $\Gr(\ell)/\Gr(\ell-1)$\\
\midrule
3 & 11 & 0.4074 & 6.1 & 1.7934 & --\\
4 & 36 & 0.4444 & 23.1 & 1.5557 & 3.273\\
5 & 113 & 0.4650 & 75.5 & 1.4968 & 3.139\\
6 & 346 & 0.4746 & 230.8 & 1.4994 & 3.062\\
7 & 1030 & 0.4710 & 681.0 & 1.5124 & 2.977\\
8 & 2995 & 0.4565 & 1967.4 & 1.5223 & 2.908\\
9 & 8580 & 0.4359 & 5605.2 & 1.5307 & 2.865\\
10 & 24303 & 0.4116 & 15816.8 & 1.5365 & 2.833\\
11 & 68283 & 0.3855 & 44324.0 & 1.5405 & 2.810\\
12 & 190698 & 0.3588 & 123572.5 & 1.5432 & 2.793\\
13 & 530152 & 0.3325 & 343156.7 & 1.5449 & 2.780\\
14 & 1468691 & 0.3071 & 949999.5 & 1.5460 & 2.770\\
\bottomrule
\end{tabular}
\end{table}

\paragraph{The root $r(\ell)$.} In~\cite{XZ}, $h(z)=(-\pi)^{\ell-1}H_\ell(\log z/\pi)\,(\log^2z+\pi^2)^{-\ell}$ and $r(\ell)=\exp(\pi\rho_\ell)$ with $\rho_\ell$ the largest root of $H_\ell$ (they note $\rho_3=\sqrt3/3$). Expanding the explicit expression for $h$ in~\cite{XZ} gives $h(z)=(-1)^{\ell-1}\operatorname{Im}\big[(\log z+i\pi)^\ell\big]\big/\big(\pi(\log^2z+\pi^2)^\ell\big)$, so its real zeros are the $\log z=\pi\cot(k\pi/\ell)$, $k=1,\dots,\ell-1$, and so $\rho_\ell=\cot(\pi/\ell)$ and $r(\ell)=\exp(\pi\cot(\pi/\ell))$. This is consistent with their table apart from the entry noted above. As $\ell\to\infty$, $\ln r(\ell)=\ell-\pi^2/(3\ell)+O(\ell^{-3})$, so $r(\ell)$ grows like $e^\ell$ and not $3^\ell$; the upper bound $A(\ell)=\exp\big(\pi\sqrt{(\ell-1)(\ell-2)/6}\big)$ for $r(\ell)$ given in~\cite{XZ} grows like $3.606^\ell$ and is far from sharp.

\paragraph{Observations (not theorems).} In Table~\ref{tab:gr}, $\Gr(\ell)/3^\ell$ peaks at $\ell=6$ and then decreases; the ratios $\Gr(\ell)/\Gr(\ell-1)$ decrease from $3.27$ at $\ell=4$ to $2.77$ at $\ell=14$ and $\Gr(\ell)/r(\ell)$ stays between $1.49$ and $1.55$. This suggests $\Gr(\ell)=\Theta(r(\ell))=\exp(\ell+O(1))$, i.e.\ growth rate $e$ and not $3$; so the upper bound $0.477\cdot3^\ell$ of Conjecture~\ref{conj:xz} should hold for large $\ell$ while its lower bound fails for all large $\ell$. A heuristic stationary-phase computation based on the integral representation of~\cite{XZ} predicts $\Gr(\ell)\sim e^{\ell+1-\gamma}$ ($\gamma$ Euler's constant); a numerical evaluation of the last sign change of the real-$n$ extension up to $\ell=100$ gives $\Gr/r$ decreasing slowly to about $1.533$, to be compared with $e^{1-\gamma}=1.526$. We have not proved this, and we did not have this evaluation reproduced independently; we include it only as a conjecture.

\section{Positivity of the reciprocal of the diagonal Pad\'e denominator of the exponential}\label{sec:pade}

\subsection{The question}
For integers $k\ge0$, $m\ge1$ let
\[
q_{k,m}(z)=\sum_{j=0}^{m}(-1)^j\binom mj\frac{(k+m-j)!}{(k+m)!}\,z^j
\]
be the denominator of the $[k/m]$ Pad\'e approximant to $e^z$ at $0$ (normalised by $q_{k,m}(0)=1$). Fasi and Higham~\cite{FH} conjectured the following (their Conjecture 3.2) in connection with a forward-error bound for the matrix exponential.

\begin{conjecture}[Fasi--Higham~\cite{FH}]\label{conj:fh}
Let $k,m\in\N$. If $m$ is odd, then all the coefficients of the series expansion of $q_{k,m}(z)^{-1}$ at $0$ are positive.
\end{conjecture}

Scardino~\cite{Scardino} proved full positivity for $k=0$ and \emph{eventual} positivity (positivity of $c_n$ for all large $n$) whenever the positive real zero of $q_{k,m}$ is the unique zero of minimal modulus, a condition he shows holds for each fixed odd $m$ and all sufficiently large $k$; he also showed that the conjecture is false in general: with $c_n(k,m)$ the coefficients of $q_{k,m}^{-1}$, he gives exact counterexamples $c_{38}(66,5)<0$, $c_{37}(44,7)<0$, $c_{42}(45,9)<0$. He writes that ``the counterexamples are not diagonal, and the diagonal case $k=m$, which is the case relevant in practice, remains open.'' Here $q_{m,m}$ is a normalised reverse Bessel polynomial. We contribute two things: a computer-assisted proof of the diagonal case for $3\le m\le361$ (Theorem~\ref{thm:pade}), and a numerical description of where positivity first fails for fixed $m$ (Section~\ref{sec:kstar}). The general diagonal statement remains open.

\begin{theorem}[computer-assisted]\label{thm:pade}
For every odd $m$ with $3\le m\le 361$, all Taylor coefficients at $0$ of $1/q_{m,m}(z)$ are positive. (For $m=1$, $1/q_{1,1}(z)=\sum_n (z/2)^n$.)
\end{theorem}

\subsection{The certificate}
Fix an odd $m$ and write $q=q_{m,m}$, $c_n$ for the coefficients of $1/q$.

\begin{lemma}\label{lem:cert}
Suppose $q$ has $m$ distinct zeros $\rho_1,\dots,\rho_m$, exactly one of which, $R=\rho_1$, is real and positive, and $|\rho_i|>R$ for $i\ge2$. Put $B_i=-1/q'(\rho_i)$, $S=\sum_{i\ge2}|B_i|$, $\rho_{\min}=\min_{i\ge2}|\rho_i|$. If $B_1>0$ and $N_0=\max\{0,\lceil\ln(S/B_1)/\ln(\rho_{\min}/R)\rceil\}$, then $c_n>0$ for all $n\ge N_0$. Consequently $c_n>0$ for all $n\ge0$ as soon as $c_0,\dots,c_{N_0-1}>0$.
\end{lemma}
\begin{proof}
Since $q(0)=1$ and $q$ has $m$ simple zeros, partial fractions give $1/q(z)=\sum_i \frac{1}{q'(\rho_i)(z-\rho_i)}$ (there is no polynomial part, the numerator having degree $0<m$), hence $c_n=\sum_iB_i\rho_i^{-n-1}$ for all $n\ge0$. For $n\ge N_0$ we have $n+1>\ln(S/B_1)/\ln(\rho_{\min}/R)$, i.e.\ $(\rho_{\min}/R)^{n+1}>S/B_1$, i.e.\ $S\rho_{\min}^{-n-1}<B_1R^{-n-1}$. Therefore
\[
c_n\ \ge\ B_1R^{-n-1}-\sum_{i\ge2}|B_i||\rho_i|^{-n-1}\ \ge\ B_1R^{-n-1}-S\rho_{\min}^{-n-1}>0 .\qedhere
\]
\end{proof}
(In the computation we take $N_0$ one larger than the displayed value to absorb rounding; this only makes the check stronger.)

\subsection{The computation}
For each odd $m$ in the range we carried out, in rigorous ball arithmetic (Arb through python-flint~\cite{Arb}):
\begin{enumerate}[leftmargin=2em]
\item isolation of the $m$ complex zeros of $q$ in disjoint balls; exactly one ball is real and positive, giving $R$; for each of the other balls a certified lower bound of the modulus exceeding $R$;
\item the residues $B_i=-1/q'(\rho_i)$ (certified balls), with $B_1>0$ certified, and the resulting bound $N_0$ of Lemma~\ref{lem:cert};
\item positivity of $c_0,\dots,c_{N_0-1}$ by inverting the power series $q$ in ball arithmetic, with the working precision raised until every one of these coefficients is a ball strictly above $0$.
\end{enumerate}
The working precision needed grows with $m$ (from 512 bits for $m\le11$ to 8192 bits for $m>300$), because the ratio $\rho_{\min}/R$ tends to $1$: it is $1.094$ for $m=3$, $1.008$ for $m=11$ and $1+3.5\cdot10^{-5}$ at $m\approx 240$ (Table~\ref{tab:pade}). The runs for $3\le m\le 361$ were recorded in log files (one line per $m$) in the repository; every one of the $180$ odd values of $m$ in this range appears with the result ``proved'', and no value is missing or failed. The largest threshold is $N_0=372{,}888$ at $m=361$ (at that point $c_n$ for $n<N_0$ are verified by series inversion with 8192 bits).

In addition, a second, independent exact certification of the first two items in the list was run: for every odd $3\le m\le361$, $\gcd(Q,Q')=1$ for the integer polynomial $Q=(2m)!\,q$ (so the zeros are simple), the zeros were isolated by Arb's certified routine for integer polynomials (which also certifies the real ones), exactly one zero is real and positive, and all other zeros have certified modulus larger than $R$.

\begin{table}[h]
\centering
\caption{Data from the certification runs (selected odd $m$). $R$ is the positive real zero of $q_{m,m}$, $\rho_{\min}/R$ the gap to the next zero in modulus, $N_0$ the threshold of Lemma~\ref{lem:cert}.}\label{tab:pade}
\begin{tabular}{rrrr}
\toprule
$m$ & $R$ & $\min_{i\ge2}|\rho_i|/R$ & $N_0$\\
\midrule
3 & 4.64437 & 1.09446 & 2\\
11 & 15.2447 & 1.00798 & 128\\
101 & 134.537 & 1.000102 & 22\,697\\
201 & 267.086 & 1.0000258 & 103\,850\\
301 & -- & -- & 251\,082\\
361 & -- & -- & 372\,888\\
\bottomrule
\end{tabular}
\end{table}

\paragraph{What is and is not proved.} Theorem~\ref{thm:pade} is a theorem for each listed $m$ individually; the proof consists of Lemma~\ref{lem:cert} together with the certified computations above, which are rigorous in the sense of ball arithmetic but are performed by a program (in the repository, \texttt{results/pade\_positivity}). It does not give the statement for all $m$. Moreover the gap $\rho_{\min}/R\to1$ and the growth of $N_0$ (the ratio $N_0/m^2$ is $2.2,\,2.6,\,2.8,\,2.9$ for $m=101,201,301,361$) indicate that this method cannot be pushed to all odd $m$; a proof for all $m$ needs a different argument.

\subsection{Where positivity first fails}\label{sec:kstar}
For fixed odd $m$ let $k^*(m)$ be the least $k$ for which a negative coefficient appears. We computed $k^*(m)$ by bisection in $k$ with interval arithmetic (each $c_n$, $n\le5k+300$, certified to have a definite sign), assuming that the set of $k$ for which negative coefficients occur is an interval $[k^*,\infty)$; we verified $k^*-1$ has no negative coefficient in that range and that $k^*,\dots,k^*+19$ all do, and checked contiguity directly only for $15\le m\le71$ and $k<400$. For small $m$ the values agree with exact rational arithmetic: $k^*(5)=66$ (first negative coefficient $n=38$), $k^*(7)=44$ ($n=37$), $k^*(9)=42$ ($n=41$), $k^*(19)=54$ ($n=68$); the first two are Scardino's examples, and $(k,m,n)=(42,9,41)$ is a smaller counterexample than his $(45,9,42)$. For $m=3$ we found no negative coefficient for any $k<1500$ (checking $n\le4k+200$). For larger $m$ the computation gives
\[
k^*(m)-2m\in\{12,13,14\}\qquad(109\le m\le 1009\ \text{sampled}),
\]
e.g.\ $k^*(109)=230$, $k^*(209)=430$, $k^*(509)=1031$, $k^*(1009)=2032$, and the first negative coefficient occurs near $n\approx k+m-3$. For $m=109$ and $m=209$ the thresholds were re-checked with exact rational arithmetic around $k^*$ (no negative coefficient among $c_n$, $n\le k+m+10$ for $k=229,429$; negative coefficients $c_{337}(230,109)$ and $c_{637}(430,209)$). In particular the diagonal $k=m$ is far from the boundary of the failure region for large $m$ ($k^*/m\to2$ numerically). This law is a numerical observation, not a theorem.

\begin{question}
Is $1/q_{m,m}$ positive for every odd $m$? Is $k^*(m)=2m+O(1)$?
\end{question}

\section{Series of products of consecutive harmonic numbers}\label{sec:harmonic}
(to be completed)

\section{Three questions of Finch on iteration}\label{sec:finch}

Finch has recently posed several questions about iterated maps~\cite{FinchDelta,FinchIII,FinchIV}. We answer three of them. The first is an elementary consequence of classical asymptotics of Fatou coordinates and is presented as such; the other two are an exact series with a numerical evaluation and a table of numerical values.

\subsection{The constant $\delta$ between two solutions of Abel's equation}

\paragraph{Setting.} Let $\theta(x)=x+\sum_{m\ge1}c_mx^{m\tau+1}$ be analytic near $0$, with an integer $\tau\ge1$ and $c_1=\gamma<0$, so that $0$ is an attracting parabolic fixed point. Put $\kappa=-\tau\gamma>0$ and
\[
a=\frac{\tau+1}2-\frac{c_2}{\gamma^2}.
\]
Let $x_0>0$ be a starting point in the basin, $x_n=\theta^n(x_0)$, so $x_n\downarrow0$. Finch considers two numerical solutions of Abel's equation $g(\theta(x))=g(x)+1$.

\begin{itemize}[leftmargin=2em]
\item The \emph{EJ solution} (\'Ecalle--Jagy). There is a unique formal series $G(x)=\kappa^{-1}x^{-\tau}+a\ln x+\sum_{m\ge1}b_mx^{m\tau}$ (no constant term) with $G(\theta(x))=G(x)+1$ (see the proof below). One sets $g_{\mathrm{EJ}}(x_0)=\lim_{n\to\infty}\big[G_K(x_n)-n\big]$ where $G_K$ is the truncation of $G$ after the term $b_Kx^{K\tau}$ ($K\ge1$; the limit exists and does not depend on $K$, see below). This is Finch's description $g(x)=\lim[g(x_n)-n]$.
\item The \emph{ML solution} (Mavecha--Laohakosol), given by the limit
\[
\tilde g(x_0)=-\tau\lim_{n\to\infty}n^{1+1/\tau}\Big[\kappa^{1/\tau}x_n-n^{-1/\tau}+\frac a{\tau^2}\,\frac{\ln n}{n^{1+1/\tau}}\Big].
\]
Finch displays this formula for the maps he studies, e.g.\ $\tilde g_{13}(x)=-2\lim n^{3/2}\big(\sqrt2x_n-n^{-1/2}+\tfrac38\ln n\,n^{-3/2}\big)$ for $\theta_{13}(x)=x(1-x^2)$ (here $\tau=2$, $\kappa=2$, $a=\tfrac32$), and the corresponding formulas for $x(1-x^k)$ and for $x(1-x)(1+\rho x)$; in all of these cases it agrees with the formula above. We take the formula above as the definition of $\tilde g$ for general $\theta$.
\end{itemize}
He conjectured exact values of $\delta=g_{\mathrm{EJ}}-\tilde g$ in many cases, in particular $\delta_{13}=-\tfrac34\ln2$, $\delta_{16}=-\tfrac23\ln3$, $\delta_{17}=-\tfrac59\ln6$, and (Addendum III of~\cite{FinchDelta}) for $\theta(x)=x(1-\sigma x^\tau-\rho x^{2\tau})$, $\sigma>0$:
\[
\delta=\Big(\frac{\tau+1}{2\tau}+\frac{\rho}{\tau\sigma^2}\Big)\ln\frac1{\tau\sigma},
\]
remarking that ``$\delta$ seems to be completely determined by the parameter $\tau$ and two coefficients $c_1$ \& $c_2$''. All these are consequences of the following statement.

\begin{proposition}\label{prop:delta}
In the setting above, the formal series $G$ exists and is unique, the limits defining $g_{\mathrm{EJ}}(x_0)$ and $\tilde g(x_0)$ exist, and
\[
\delta=g_{\mathrm{EJ}}(x_0)-\tilde g(x_0)=-\frac a\tau\ln\kappa=\Big(\frac{\tau+1}{2\tau}-\frac{c_2}{\tau\gamma^2}\Big)\ln\frac1{\tau|\gamma|}\qquad\text{for every }x_0.
\]
For $\theta=x(1-\sigma x^\tau-\rho x^{2\tau})$ ($c_1=-\sigma$, $c_2=-\rho$) this is exactly Finch's Addendum III formula.
\end{proposition}

\begin{remark}
The logarithmic coefficient $a=\tfrac{\tau+1}2-c_2/c_1^2$ is classical: it is (up to normalisation) \'Ecalle's \emph{r\'esidu it\'eratif} of the parabolic germ, which is the coefficient of the logarithm in the Fatou coordinate (see for example Milnor's survey and Lanford--Yampolsky~\cite{LY}). Given this, Proposition~\ref{prop:delta} is a short corollary of the definitions of the two algorithms, and we do not claim the asymptotics as new; what we supply is a complete derivation of Finch's conjectured constants.
\end{remark}

\begin{proof}
\emph{The formal series.} Put $\eps=x^\tau$. Then
\[
\theta(x)^{-\tau}=x^{-\tau}\big(1+\gamma\eps+c_2\eps^2+O(\eps^3)\big)^{-\tau}=\eps^{-1}-\tau\gamma+\Big(\tfrac{\tau(\tau+1)}2\gamma^2-\tau c_2\Big)\eps+O(\eps^2),
\]
\[
\ln\frac{\theta(x)}{x}=\gamma\eps+O(\eps^2),\qquad \theta(x)^{m\tau}-x^{m\tau}=m\tau\gamma\,\eps^{m+1}+O(\eps^{m+2}).
\]
For $G=Ax^{-\tau}+a'\ln x+\sum_{m\ge1}b_mx^{m\tau}$ we get $G(\theta(x))-G(x)=-A\tau\gamma+\big[A(\tfrac{\tau(\tau+1)}2\gamma^2-\tau c_2)+a'\gamma\big]\eps+\sum_{j\ge2}(\cdots)\eps^j$, where the coefficient of $\eps^j$ for $j\ge2$ contains $b_{j-1}$ with the nonzero factor $(j-1)\tau\gamma$ plus terms involving only $A,a',b_1,\dots,b_{j-2}$. Equating with $1$ gives successively $A=-1/(\tau\gamma)=1/\kappa$; then $a'=-A(\tfrac{\tau(\tau+1)}2\gamma^2-\tau c_2)/\gamma=\tfrac{\tau+1}2-\tfrac{c_2}{\gamma^2}=a$ (using $A=-1/(\tau\gamma)$); then $b_1,b_2,\dots$ uniquely. So $G$ exists, is unique, and the truncation satisfies $G_K(\theta(x))-G_K(x)-1=O(x^{(K+1)\tau})$.

\emph{Existence of $g_{\mathrm{EJ}}$.} Since $x_n\to0$ monotonically and the recurrence gives $x_{n+1}/x_n=1+\gamma x_n^\tau+O(x_n^{2\tau})$, we have $x_n^{-\tau}=\kappa n(1+o(1))$ (indeed $x_{n+1}^{-\tau}-x_n^{-\tau}=\kappa+O(x_n^\tau)\to\kappa$, and Ces\`aro's lemma applies), so $x_j^{(K+1)\tau}=O(j^{-(K+1)})$ is summable for $K\ge1$. Hence $G_K(x_n)-n-G_K(x_0)=\sum_{j<n}[G_K(x_{j+1})-G_K(x_j)-1]$ converges, so $g_{\mathrm{EJ}}(x_0)$ exists; it does not depend on $K\ge1$ because $b_{K+1}x_n^{(K+1)\tau}\to0$. Moreover, since $b_mx_n^{m\tau}\to0$,
\begin{equation}\label{eq:orbit}
\kappa^{-1}x_n^{-\tau}+a\ln x_n=n+g_{\mathrm{EJ}}(x_0)+o(1).
\end{equation}

\emph{Asymptotics of the orbit.} From \eqref{eq:orbit}, $x_n^{-\tau}=\kappa n+\kappa g_{\mathrm{EJ}}-\kappa a\ln x_n+o(1)$. Since $x_n\to0$ and $\ln x_n=O(\ln x_n^{-\tau})$ we first get $x_n^{-\tau}=\kappa n+O(\ln n)$, hence $\ln x_n=-\tfrac1\tau\ln(\kappa n)+o(1)$, and then
\[
x_n^{-\tau}=\kappa n\,(1+\eps_n),\qquad \eps_n=\frac{\tfrac a\tau\ln(\kappa n)+g_{\mathrm{EJ}}(x_0)}{n}+o\Big(\frac1n\Big)=O\Big(\frac{\ln n}n\Big).
\]
Therefore $\kappa^{1/\tau}x_n=n^{-1/\tau}(1+\eps_n)^{-1/\tau}=n^{-1/\tau}\big[1-\tfrac1\tau\eps_n+O(\ln^2n/n^2)\big]$ and
\[
\kappa^{1/\tau}x_n-n^{-1/\tau}+\frac a{\tau^2}\frac{\ln n}{n^{1+1/\tau}}=-\frac1\tau\,n^{-1-1/\tau}\Big[\frac a\tau\ln\kappa+g_{\mathrm{EJ}}(x_0)\Big]+o\big(n^{-1-1/\tau}\big),
\]
where the $\ln n$ terms cancelled. Multiplying by $-\tau n^{1+1/\tau}$ shows that the limit defining $\tilde g(x_0)$ exists and equals $g_{\mathrm{EJ}}(x_0)+\tfrac a\tau\ln\kappa$. \end{proof}

\begin{remark}[Checks]
The proposition reproduces every value in Finch's paper that we tested: $\delta_{13}$, $\delta_{14}=0$, $\delta_{15}$, $\delta_{16}$, $\delta_{17}$, $\delta(k)=-\frac{k+1}{2k}\ln k$ for $x(1-x^k)$ and for $x(1-x)^k$, $\delta(\rho)$, and the transcendental examples $\ln(1+x)$ ($-\tfrac13\ln2$), $xe^{-x}$ ($0$), $x/(1+x^2)$ ($-\tfrac14\ln 2$), $\sin x$ ($\tfrac35\ln3$) and the Fresnel cosine ($\tfrac{55}{108}\ln\tfrac{10}{\pi^2}$). Numerically, Finch's printed 100-digit values $g_{13}(\tfrac12)=1.18935\ldots$ and $\tilde g_{13}(\tfrac12)=1.70921\ldots$ differ by $-0.5198604\ldots=-\tfrac34\ln2$; and a fresh implementation at 40 digits with $4\cdot10^5$ orbit steps reproduces $-0.5198604$ to seven digits (the ML limit converges like $\ln^2n/n$ and needs extrapolation). These checks use his definitions of the algorithms; we have used his version v6, as an earlier draft (v3) contained a programming error acknowledged in his paper.
\end{remark}

\subsection{The second derivative $C''(2)$ of an orbit limit}
Finch~\cite{FinchIII} studies $x_{n+1}=x_n+1+1/x_n^2$ and the limit $C(x_0)=\lim_{n\to\infty}(x_n-n)$. He computes by finite differences $C(2)=2.59878685582487\ldots$, $C'(2)=0.66156132404868\ldots$, $C''(2)\approx0.37462642198301734111$, notes that an earlier series estimate $0.3909\ldots$ disagrees, and writes: ``An independent confirmation of the $0.3746$ estimate remains open.''

Let $x_k$ be the orbit of $x_0$, $f(x)=x+1+x^{-2}$ so $f'(x)=1-2x^{-3}$. Since $x_k-k\to C$ and $\prod_k f'(x_k)$ converges ($x_k\sim k$, $\sum x_k^{-3}<\infty$), $C'(x_0)=\prod_{k\ge0}(1-2x_k^{-3})$. Writing $D_k=\prod_{j<k}(1-2x_j^{-3})=\dd x_k/\dd x_0$ and differentiating once more, $\dd^2x_k/\dd x_0^2=:E_k$ satisfies $E_{k+1}=f'(x_k)E_k+f''(x_k)D_k^2$, with $f''(x)=6x^{-4}$, and $C''(x_0)=\lim_kE_k$. Solving the linear recursion,
\begin{equation}\label{eq:Cpp}
C''(x_0)=C'(x_0)\sum_{j\ge0}\frac{6\,D_j}{x_j^{4}\,(1-2x_j^{-3})}=C'(x_0)\sum_{j\ge0}\frac{6\,D_j}{x_j(x_j^3-2)} .
\end{equation}
(Finch's series in~\cite{FinchIII}, $C'(x_0)\sum_n\big(1-2x_n^{-3}\big)^{-1}6x_n^{-4}$ under a uniformity assumption, lacks the factor $D_j$; this is why it gave $0.3909$.) The interchange of the limit $k\to\infty$ with differentiation is justified by the uniform summability of the terms; we have not written out the details, and \eqref{eq:Cpp} converges only like $O(1/N)$ in the number $N$ of terms. For $x_0=2$ we obtained, with the asymptotic expansion $x_n-n=C-\frac1n+\frac{C-1/2}{n^2}+\cdots$ (all logarithmic terms vanish) and exact orbit recursions for the derivatives,
\[
C''(2)=0.374626421983017341116776\ldots,
\]
in agreement with Finch's $0.37462642198301734111$ to all 20 digits he printed. A second, independent implementation (fresh code; $N=10^4$ and $3\cdot10^4$ orbit steps agree to 33--35 digits) gives $0.374626421983017341116776248793$, agreeing with the digits above, and a finite-difference evaluation at high precision agrees to 22 digits; the series \eqref{eq:Cpp} summed over $2\cdot10^5$ terms gives $0.3746264219924$, consistent with the expected tail $O(1/N)$. None of these computations is interval-certified.

\subsection{Half-iterates of $b^x$ for $b=4/3,\,5/4,\,6/5$}
Finch~\cite{FinchIV} writes: ``It would be good to see the values of $T(4/3,1/2)$, $T(5/4,1/2)$ and $T(6/5,1/2)$ someday.'' Here $T(b,x)$ is tetration, and $T(b,\tfrac12)=f^{[1/2]}(1)$ for $f(x)=b^x$, where for $b<e^{1/e}$ we use the regular (Koenigs--Schr\"oder) iteration at the attracting fixed point $L$ ($b^L=L$, $L<e$): with multiplier $\lambda=f'(L)=L\ln b\in(0,1)$ and Koenigs function $\phi(f(x))=\lambda\phi(x)$, $\phi'(L)=1$, one puts $f^{[1/2]}=\phi^{-1}(\sqrt\lambda\,\phi(\cdot))$. This is the real-analytic half-iterate that is analytic at $L$; other solutions of $g\circ g=f$ exist and are not analytic at $L$.

\begin{table}[h]
\centering
\caption{Half-iterates $T(b,\tfrac12)=f^{[1/2]}(1)$, $f(x)=b^x$ (regular iteration at the attracting fixed point $L$). Two independent methods (power series of the inverse Koenigs function; limit formulas at 500-digit working precision) agree to the digits shown.}\label{tab:tet}
\begin{tabular}{llll}
\toprule
$b$ & $L$ & $\lambda$ & $T(b,\tfrac12)$\\
\midrule
$4/3$ & 1.57167727089176302 & 0.45214337 & 1.205837434043095594941839306507112566900\\
$5/4$ & 1.35220351364068086 & 0.30173549 & 1.163595361965008849025943273358436312883\\
$6/5$ & 1.25773454137652642 & 0.22931212 & 1.136262486727128084185300918647426025589\\
\bottomrule
\end{tabular}
\end{table}

As a check of the convention, the same construction for $b=\sqrt2$ (attracting fixed point $L=2$, $\lambda=\ln2$) reproduces the value $f^{[1/2]}(1)=1.24362162766852180429509898\ldots$ printed in~\cite{FinchIV}, Section 1.5 (in a quick check with 120 iterations, to about 20 digits). The values are numerical, not interval-certified: the two methods agree to at least 50 digits, the functional equation $g(g(x))=b^x$ holds to better than $10^{-85}$ at several points, and the digits shown are well inside the observed convergence.

\section{Gurland's ratio of gamma functions}\label{sec:gurland}

Wi\'sniewska~\cite{Wis} studies the modified Gurland ratio
\[
G^\star(x,y)=\frac{\Gamma(1+x)\,\Gamma(1+y)}{\Gamma^2\!\big(1+\frac{x+y}2\big)},\qquad x,y>0,
\]
and proves, for integers $m\ge2$, the finite expansion
\[
\ln G^\star(x,y)=S_m(x,y)+R_{nm}(x,y),\qquad S_m(x,y)=\sum_{k=1}^{m-1}\frac1k\Big(\frac{x-y}2\Big)^{2k}\zeta\Big(2k,1+\frac{x+y}2\Big),
\]
with $0\le R_{nm}(x,y)\le\varepsilon_n^m(x,y)$, where $\zeta(s,a)=\sum_{n\ge0}(n+a)^{-s}$ is the Hurwitz zeta function. From this she deduces $\ln G^\star=S_\infty:=\lim_mS_m$ under the condition $|x-y|<2(1+\sqrt{xy})$ and ends Section 2 with three open problems: (1) ``Does $R_{nm}(x,y)\to0$ hold for every $x,y>0$?'' (2) ``Does $\ln G^\star(x,y)=S_\infty(x,y)$ hold generally, and if not, how large the difference can be?'' (3) ``How is $t(x,y)$ located in the interval $[\sqrt{xy},\frac{x+y}2]$?'' where $t$ is defined by $\ln G^\star=\big(\frac{x-y}2\big)^2\zeta(2,1+t)$ (her Corollary 3).

\subsection{Problems 1 and 2}
\begin{proposition}\label{prop:gurland}
For all $x,y>0$, with $h=\frac{x+y}2$, $d=\frac{x-y}2$,
\[
\ln G^\star(x,y)=\sum_{k=1}^\infty\frac1k\,d^{2k}\,\zeta(2k,1+h),
\]
the series converging absolutely, with radius of convergence exactly $1+h$ in the variable $d$; moreover $|d|\le h+1-(1+\min(x,y))$ for all $x,y>0$, so the convergence is geometric with ratio at most $(d/(1+h))^2<1$. In particular $R_{nm}(x,y)\to0$ as $m\to\infty$ and $\ln G^\star=S_\infty$ for every $x,y>0$.
\end{proposition}
\begin{proof}
Put $a=1+h$ and $f(z)=\ln\Gamma(a+z)$ for $|z|<a$. Here $\ln\Gamma$ is the branch that is real on $(0,\infty)$; it is holomorphic on the half-plane $\operatorname{Re}w>0$ (the zeros of $1/\Gamma$ are absent and the poles of $\Gamma$ lie on the non-positive real axis). For $|z|<a$ we have $\operatorname{Re}(a+z)>0$, so $f$ is holomorphic in the disc $|z|<a$, and its Taylor series at $0$ converges to $f$ there. Its coefficients are $f^{(j)}(0)=\psi^{(j-1)}(a)$ for $j\ge1$, and $\psi^{(n)}(a)=(-1)^{n+1}n!\,\zeta(n+1,a)$ for $n\ge1$; hence $f^{(2k)}(0)=(2k-1)!\,\zeta(2k,a)$ for $k\ge1$. Since $1+x=a+d$ and $1+y=a-d$, we have $\ln G^\star=f(d)+f(-d)-2f(0)$, and $|d|<a$ holds for all $x,y>0$ because $a-|d|=1+\min(x,y)>1$. Adding the Taylor series at $d$ and $-d$, the odd terms cancel:
\[
\ln G^\star=2\sum_{k\ge1}\frac{f^{(2k)}(0)}{(2k)!}d^{2k}=\sum_{k\ge1}\frac1k\,\zeta(2k,a)\,d^{2k}.
\]
The terms are nonnegative, and $\zeta(2k,a)\ge a^{-2k}$, so for $|d|=a$ the series would diverge (it is bounded below by $\sum\frac1k$); thus the radius is exactly $a$. For $|d|<a$, $\zeta(2k,a)=a^{-2k}(1+O((a/(a+1))^{2k}))$ and the terms are $O\big(k^{-1}(d/a)^{2k}\big)$. Finally $S_m$ is the $(m-1)$-st partial sum, so $R_{nm}=\ln G^\star-S_m\to0$.
\end{proof}

Thus the answer to Problems 1 and 2 is ``yes'' and the difference $\ln G^\star-S_\infty$ is $0$. The argument is elementary (it is just the radius of convergence of the Taylor series of $\ln\Gamma$ at $1+h$), and we do not regard it as deep; it was not stated in~\cite{Wis}, whose sufficient condition $|x-y|<2(1+\sqrt{xy})$ is therefore not necessary. As a sanity check, for $(x,y)=(0.01,100)$, where $|d|/a=0.98$, the series evaluated with 60--120 working digits agrees with $\ln G^\star$ and the remainders behave like the predicted $\sim r^{2N}/(N(1-r^2))$, $r=|d|/a$ (about $10^{-10}$ after 500 terms).

\subsection{Problem 3: partial information (numerical)}
Write $g=\sqrt{xy}$, so $g\le t\le h$. From Proposition~\ref{prop:gurland}, as $d\to0$ with $h$ fixed,
\[
\frac{\ln G^\star}{d^2}=\psi'(1+h)+\frac{d^2}{12}\psi'''(1+h)+O(d^4),\qquad\text{so}\qquad t=h+\frac{d^2}{12}\frac{\psi'''(1+h)}{\psi''(1+h)}+O(d^4),
\]
while $h-g=d^2/(2h)+O(d^4)$. Hence the relative position $\vartheta=(h-t)/(h-g)$ of $t$ between $h$ and $g$ satisfies $\vartheta\to\frac h6\,\frac{\psi'''(1+h)}{-\psi''(1+h)}$ as $d\to0$, which tends to $\frac13$ as $h\to\infty$ (using $\psi''(z)\sim-z^{-2}$, $\psi'''(z)\sim2z^{-3}$). So no bound of the form $t\ge h-c\,(h-g)$ can hold with $c<\frac13$. In a numerical search (20\,000 random points with $x,y\in[10^{-3},10^{3.5}]$, the largest 200 values re-evaluated with 40-digit arithmetic) we found $\vartheta<\frac13$ with maximum $0.33327$, which suggests:

\begin{conjecture}\label{conj:gurland3}
For all $x\ne y$ in $(0,\infty)$, $t(x,y)>\dfrac{x+y+\sqrt{xy}}3$ (equivalently $\vartheta<\tfrac13$). The constant $\tfrac13$ cannot be improved.
\end{conjecture}
We have no proof of this inequality. (It is a statement about a Hurwitz-zeta quantity for which a termwise comparison in the Laplace variable fails.)

\section{Completely monotonic degrees of Stirling remainders}\label{sec:stirling}

\subsection{The question}
A function $f\in C^\infty(0,\infty)$ is completely monotonic (CM) if $(-1)^jf^{(j)}\ge0$ for all $j\ge0$. For $f\to0$ at infinity, the \emph{real completely monotonic degree} is $d_{CM}(f)=\sup\{\rho\ge0: x^\rho f(x)\ \text{is CM}\}$. For $n\ge0$ let
\[
R_n(x)=(-1)^n\Big[\ln\Gamma(x)-\big(x-\tfrac12\big)\ln x+x-\tfrac12\ln2\pi-\sum_{k=1}^n\frac{B_{2k}}{2k(2k-1)x^{2k-1}}\Big],\qquad x>0,
\]
the signed Stirling remainders, with Bernoulli numbers $B_j$. Krasniqi~\cite{Kr} proves
\[
\tfrac83<d_{CM}(R_2)<3,\qquad \tfrac{32}9<d_{CM}(-R_2')<4
\]
(his Theorem 1.2), and, for $\delta_{n}:=d_{CM}(R_n)-(2n-2)$, that $\limsup_n(2n-1)\sup_m\delta_{n,m}\le 2\pi^2/J$ with $J=\int_0^{2\pi}\frac{1-\cos u}u\,\dd u$ (Theorem 1.3; here $\delta_{n,m}$ is the analogous excess for the $m$-th derivative, and $\delta_{n,0}=\delta_n$). He adds that ``the exact real degrees and the optimal asymptotic constants remain open.'' We computed the degrees numerically. The results below are numerical, not proved.

\subsection{Method}
For $n\ge2$, $R_n(x)=\int_0^\infty e^{-xt}h_n(t)\,\dd t$ with $h_n(t)=(-1)^nt^{-1}\big[\frac1{e^t-1}-\frac1t+\frac12-\sum_{k\le n}\frac{B_{2k}t^{2k-1}}{(2k)!}\big]\ge0$, and $-R_n'$ has density $t\,h_n(t)$. By Bernstein's theorem and uniqueness of the Laplace transform, $x^\rho f$ is CM if and only if the Riemann--Liouville derivative $D^\rho h$ of the density is nonnegative on $(0,\infty)$ (for $\rho$ below the order of vanishing of $h$ at $0$ there are no boundary terms). We computed $\rho_c=\sup\{\rho: D^\rho h\ge0\}$; at $\rho_c$ the function $D^\rho h$ first touches zero at a double root $t_*\approx11.47$ (for $R_2$) and $\approx8.04$ (for $-R_2'$). Two independent implementations were used (one with Chebyshev interpolation and QUADPACK algebraic-weight quadrature, one with Gauss--Legendre quadrature after a change of variables), and the second implementation additionally checked the result \emph{without} the Laplace representation by computing exact derivatives of $x^\rho f$ (Leibniz rule and polygamma functions at 300 digits) up to order $200$: all derivatives of the right sign for $\rho=\rho_c-0.02$ and $\rho_c-0.003$, and violations for $\rho=\rho_c+0.02$.

\subsection{Results}
\begin{proposition}[numerical]\label{prop:stirling}
To the digits shown (agreement of the two independent implementations to at least 12 digits),
\[
d_{CM}(R_2)=2.8667042552\ldots,\qquad d_{CM}(-R_2')=3.7984354776\ldots,
\]
both inside the intervals proved in~\cite{Kr}.
\end{proposition}
Table~\ref{tab:stirling} gives $d_{CM}(R_n)$ for further $n$ and the normalised excess $(2n-1)\delta_n$. These values come from our first implementation; the entries marked $^\dagger$ were also recomputed by the independent implementation.

\begin{table}[h]
\centering
\caption{$d_{CM}(R_n)$ and $(2n-1)\delta_n$, $\delta_n=d_{CM}(R_n)-(2n-2)$; and $d_{CM}(-R_n')$ for small $n$. ($^\dagger$: recomputed independently.)}\label{tab:stirling}
\begin{tabular}{rllr}
\toprule
$n$ & $d_{CM}(R_n)$ & $(2n-1)\delta_n$ & $d_{CM}(-R_n')$\\
\midrule
2 & 2.8667042552$^\dagger$ & 2.6001 & 3.7984354776$^\dagger$\\
3 & 4.7074638766$^\dagger$ & 3.5373 & 5.6534285885$^\dagger$\\
4 & 6.5993536696 & 4.1955 & 7.5587377578\\
5 & 8.5213225063$^\dagger$ & 4.6919 & 9.4900039902$^\dagger$\\
6 & 10.4617284606 & 5.0790 & 11.4369053276\\
8 & 14.3760290864 & 5.6404 & 15.3593656569\\
10 & 18.3171745147$^\dagger$ & 6.0263 & --\\
14 & 26.2415343007 & 6.5214 & --\\
18 & 34.1950137509 & 6.8255 & --\\
22 & 42.1635160400$^\dagger$ & 7.0312 & --\\
27 & 52.1360471128 & 7.2105 & --\\
\bottomrule
\end{tabular}
\end{table}

\paragraph{Numerical observation on the sharpness of Krasniqi's constant.} The numbers $(2n-1)\delta_n$ increase monotonically with $n$ in the table. Extrapolating the sequence for $n=2,\dots,27$ by least-squares polynomials of degree $4$--$6$ in $1/n$ or $1/(n+\frac12)$ (using the data for $n\ge8,10,12$) gives limits between $8.0967$ and $8.0977$, median $8.0975$, to be compared with Krasniqi's bound $2\pi^2/J=2\pi^2/\mathrm{Cin}(2\pi)=8.097627\ldots$. The data are therefore compatible with $\lim_n(2n-1)\delta_n=2\pi^2/J$, i.e.\ with his upper bound being attained (at $m=0$); the extrapolation is not a proof, and the spread above is a rough indication of its uncertainty (no error bound).

\paragraph{Caveats.} The computations are quadrature computations with the stability checks described, not interval arithmetic; non-negativity of $D^\rho h$ below $\rho_c$ was checked on grids, and the root $\rho_c$ was located as a tangency (double root) of $D^\rho h$. The entries marked $^\dagger$ agree with the independent implementation to at least 10 digits (for $n=22$ the independent value needed many quadrature nodes to converge; its quoted uncertainty is $3\cdot10^{-11}$); the other entries come from one implementation only.

\section{Open problems}\label{sec:open}

We collect questions that our work leaves open. Items 1--3 are conjectures we could neither prove nor refute; the others are open questions from the cited papers on which we have only partial information.

\begin{enumerate}[leftmargin=2em]
\item \textbf{Diagonal Pad\'e positivity.} Is $1/q_{m,m}(z)$ positive (all Taylor coefficients) for every odd $m$? Theorem~\ref{thm:pade} covers $3\le m\le361$. Is $k^*(m)=2m+O(1)$ (Section~\ref{sec:kstar})? The method of Section~\ref{sec:pade} cannot reach all $m$ because the gap $\rho_{\min}/R\to1$.
\item \textbf{Growth of $\Gr(\ell)$.} Prove that $\log\Gr(\ell)=\ell+O(1)$, or the sharper observed law $\Gr(\ell)\sim e^{\ell+1-\gamma}$ (Section~\ref{sec:gregory}). In particular decide whether the upper bound $\Gr(\ell)<0.477\cdot3^\ell$ of~\cite{XZ} holds for all $\ell\ge5$ (it holds for $5\le\ell\le14$).
\item \textbf{Gurland's inequality.} Is $t(x,y)>(x+y+\sqrt{xy})/3$ for all $x\ne y$ (Conjecture~\ref{conj:gurland3})? This would answer the third open problem of~\cite{Wis} sharply.
\item \textbf{Nonnegativity for the Painlev\'e-type expansion.} For $u''u-(u')^2=e^{-x}P(u)-1$ with $P(0)=0$, $P'(0)=\tfrac12P''(0)$ and all coefficients of $P$ nonnegative, are all coefficients of the polynomials $p_k$ in $u=\sum_kp_k(x)e^{-kx}$ nonnegative (Conjecture 7.2 of~\cite{HH}), and does the series converge (Conjecture 7.3 there)? The existence of the polynomials $p_k$ is proved in~\cite{HH2}. Our exact rational computations (not independently reproduced) found nonnegative coefficients in all cases tried: $P=t^3$ up to $k=40$, $P=t^4$ up to $k=30$, $P=t^6$ up to $k=20$.
\item \textbf{Sharpness in Krasniqi's asymptotic bound.} Is $\lim_{n\to\infty}(2n-1)\,(d_{CM}(R_n)-(2n-2))=2\pi^2/\mathrm{Cin}(2\pi)$ (Section~\ref{sec:stirling})? The numerical evidence suggests that the bound of~\cite{Kr} is attained.
\item \textbf{Harmonic products.} Does the top-weight law $\text{weight-}(q+2)\text{ part of }S_q=\frac1{q!}\sum_{n\ge1}H_n/n^{q+1}$ hold for all $q$ (checked exactly for $q\le9$; Section~\ref{sec:harmonic})? Is $\zeta(2)\zeta(3)$ in the $\Q$-span of $\zeta(2),\dots,\zeta(5)$ (this is equivalent to the conjecture of Furdui and Sîntămărian holding for $q=3$, see~\cite{CN})?
\item \textbf{Uniqueness in Fan's conjecture.} Is the equilibrium measure of the vertex set of the regular $N$-gon, $N$ odd, $r>2$, unique up to rotation? Theorem~\ref{thm:riesz} gives existence of an equilibrium measure with the stated support and masses.
\end{enumerate}

\begin{thebibliography}{99}
\bibitem{Fan} Q.~Fan, \emph{Riesz capacity ratios with negative exponents}, arXiv:2609.11186 (10 Sep 2026).
\bibitem{XZ} C.~Xu and J.~Zhao, \emph{Rational approximations for reciprocals of multiple zeta values and trivariate Cauchy numbers}, arXiv:2609.11072 (10 Sep 2026).
\bibitem{FH} M.~Fasi and N.~J.~Higham, \emph{An arbitrary precision scaling and squaring algorithm for the matrix exponential}, MIMS EPrint 2018.36, Manchester Institute for Mathematical Sciences, version of 19 March 2019, \url{https://eprints.maths.manchester.ac.uk/2691/}.
\bibitem{Scardino} F.~Scardino, \emph{On a positivity conjecture of Fasi and Higham for Pad\'e approximants to the matrix exponential}, arXiv:2610.39051 (28 Oct 2026).
\bibitem{CN} W.~J.~Chen and V.~Nguyen, \emph{A series involving a product of four consecutive harmonic numbers}, arXiv:2507.19502 (10 Jul 2025).
\bibitem{FS} P.~Flajolet and B.~Salvy, \emph{Euler sums and contour integral representations}, Experiment.\ Math.\ \textbf{7} (1998), no.~1, 15--35 (bibliographic details from the Project Euclid listing).
\bibitem{XW} C.~Xu and W.~Wang, \emph{Explicit formulas of Euler sums via multiple zeta values}, arXiv:1805.08056 (2018).
\bibitem{FinchDelta} S.~Finch, \emph{Half-iterates and delta conjectures}, arXiv:2506.07625 (version v6, 3 Sep 2025).
\bibitem{FinchIII} S.~Finch, \emph{Exercises in iterational asymptotics III}, arXiv:2503.13378 (17 Mar 2025).
\bibitem{FinchIV} S.~Finch, \emph{Exercises in iterational asymptotics IV}, arXiv:2509.24918 (29 Sep 2025).
\bibitem{FinchAB} S.~Finch, \emph{Abel's functional equation and interrelations}, arXiv:2503.00579 (version v3, 18 Mar 2025).
\bibitem{LY} O.~Lanford III and M.~Yampolsky, \emph{The fixed point of the parabolic renormalization operator}, arXiv:1108.2801 (2011), Section~2.2 (Proposition~2.14).
\bibitem{Wis} H.~Wi\'sniewska, \emph{Some inequalities for Gurland's ratio of the gamma functions}, arXiv:2512.07028 (7 Dec 2025).
\bibitem{Kr} V.~Krasniqi, \emph{Noninteger completely monotonic degrees of Stirling remainders and their derivatives}, arXiv:2610.33699 (24 Oct 2026).
\bibitem{HH} R.~Habib and R.~Hildebrand, \emph{Necessary conditions for the existence of exponential-polynomial expansions for solutions of certain differential equations}, arXiv:2606.21145 (19 Jun 2026).
\bibitem{HH2} R.~Hildebrand and R.~Habib, \emph{Sufficient conditions for the existence of exponential-polynomial expansions for solutions of certain differential equations}, arXiv:2606.24182 (23 Jun 2026).
\bibitem{CKL} S.~Chavan, M.~Kobayashi and J.~Layja, \emph{Integral evaluation of odd Euler sums, multiple $t$-value $t(3,2,\ldots,2)$ and multiple zeta value $\zeta(3,2,\ldots,2)$}, arXiv:2111.07097 (13 Nov 2021).
\bibitem{Charlton} S.~Charlton, \emph{On motivic multiple $t$ values, Saha's basis conjecture, and generators of alternating MZV's}, arXiv:2112.14613 (29 Dec 2021).
\bibitem{Zheng} H.~Zheng, \emph{Samart's conjecture $n_4(81)=40M_7$: the exact CM evaluation and the two obstructions. A status report}, arXiv:2608.02265 (3 Aug 2026).
\bibitem{CL} C.~Clark and R.~S.~Laugesen, \emph{Maximizing Riesz capacity ratios: conjectures and theorems}, SIAM J.\ Math.\ Anal.\ \textbf{57} (2025), no.~5, 5683--5717 (as listed in~\cite{Fan}; not read by us).
\bibitem{Bj} G.~Bj\"orck, \emph{Distributions of positive mass, which maximize a certain generalized energy integral}, Ark.\ Mat.\ \textbf{3} (1956), 255--269 (as listed in~\cite{Fan}; not read by us).
\bibitem{Arb} F.~Johansson and contributors, \emph{Arb}, a C library for arbitrary-precision ball arithmetic (since 2023 part of FLINT), \url{https://arblib.org}; used through the python-flint bindings.
\end{thebibliography}

\end{document}
